% Marshalling and unmarshalling the arguments of sum(count, values).
% The pointer value is not sent; the pointed-to array is copied into the
% buffer and recreated at a different address on the server.
% Usage: % Marshalling and unmarshalling the arguments of sum(count, values).
% The pointer value is not sent; the pointed-to array is copied into the
% buffer and recreated at a different address on the server.
% Usage: % Marshalling and unmarshalling the arguments of sum(count, values).
% The pointer value is not sent; the pointed-to array is copied into the
% buffer and recreated at a different address on the server.
% Usage: % Marshalling and unmarshalling the arguments of sum(count, values).
% The pointer value is not sent; the pointed-to array is copied into the
% buffer and recreated at a different address on the server.
% Usage: \input{figures/l13-marshal}   (width ~117 mm, fits the 120 mm body)
\begin{tikzpicture}[x=1mm, y=1mm, font=\footnotesize\sffamily,
  >={Stealth[length=1.6mm]},
  cell/.style={draw, fill=white, minimum height=5mm, inner sep=0pt,
               font=\scriptsize\sffamily},
  var/.style={cell, minimum width=16mm},
  elem/.style={cell, minimum width=7mm},
  buf/.style={cell, minimum width=7mm, fill=ln-accent!15},
  vname/.style={font=\scriptsize\sffamily, anchor=east},
  hdr/.style={font=\scriptsize\sffamily\bfseries, text=ln-muted, anchor=base},
  lab/.style={font=\scriptsize\sffamily, text=ln-muted},
  addr/.style={font=\scriptsize\ttfamily, text=ln-muted, anchor=south west,
               inner sep=0.5pt},
  op/.style={font=\scriptsize\sffamily, anchor=south}]
  % ---------------- client address space
  \node[hdr] at (15,42) {\iflangja{クライアントのアドレス空間}{client address space}};
  \node[lab, anchor=west] at (0,36.5) {\iflangja{スタック}{stack}};
  \node[var] (cc) at (21,31) {3};
  \node[var] (cv) at (21,25) {\texttt{0x5a10}};
  \node[vname] at (12.5,31) {\texttt{count}};
  \node[vname] at (12.5,25) {\texttt{values}};
  \node[elem] (c1) at (11.5,9) {7};
  \node[elem] (c2) at (18.5,9) {$-2$};
  \node[elem] (c3) at (25.5,9) {300};
  \node[addr, anchor=north west] at (8,6.3) {0x5a10};
  \node[lab, anchor=east] at (7.6,9) {\iflangja{ヒープ}{heap}};
  \draw[->] ([xshift=-4mm]cv.south) -- (c1.north);
  % ---------------- message buffer
  \node[hdr] at (60,42) {\iflangja{メッセージバッファ}{message buffer}};
  \foreach \v [count=\i from 0] in {3,7,$-2$,300} {
    \node[buf] (b\i) at (49.5+7*\i,20) {\v};
    \node[font=\tiny\sffamily, text=ln-muted] at (46+7*\i,24.5) {\the\numexpr4*\i\relax};
  }
  \node[font=\tiny\sffamily, text=ln-muted] at (74,24.5) {16};
  \draw[ln-muted] (46.4,16.8) -- (46.4,15.8) -- (52.6,15.8) -- (52.6,16.8);
  \node[lab, anchor=north] at (49.5,15.3) {\texttt{count}};
  \draw[ln-muted] (53.4,16.8) -- (53.4,15.8) -- (73.6,15.8) -- (73.6,16.8);
  \node[lab, anchor=north] at (63.5,15.3)
    {\iflangja{要素}{elements}};
  % ---------------- server address space
  \node[hdr] at (103,42) {\iflangja{サーバのアドレス空間}{server address space}};
  \node[lab, anchor=west] at (88,36.5) {\iflangja{スタック}{stack}};
  \node[var] (sc) at (109,31) {3};
  \node[var] (sv) at (109,25) {\texttt{0x9c40}};
  \node[vname] at (100.5,31) {\texttt{count}};
  \node[vname] at (100.5,25) {\texttt{values}};
  \node[elem] (s1) at (99.5,9) {7};
  \node[elem] (s2) at (106.5,9) {$-2$};
  \node[elem] (s3) at (113.5,9) {300};
  \node[addr, anchor=north west] at (96,6.3) {0x9c40};
  \node[lab, anchor=east] at (95.6,9) {\iflangja{ヒープ}{heap}};
  \draw[->] ([xshift=-4mm]sv.south) -- (s1.north);
  % ---------------- operations
  % Japanese: below the arrows; above them the wider "アンマーシャル"
  % runs into the buffer offset "16".
  \draw[->, thick, ln-accent] (31,20) -- (44.5,20);
  \node[op, anchor=\iflangja{north}{south}] at (37.8,\iflangja{19.2}{20.8}) {\iflangja{マーシャル}{marshal}};
  \draw[->, thick, ln-accent] (75.5,20) -- (87,20);
  \node[op, anchor=\iflangja{north}{south}] at (\iflangja{82.3}{81.3},\iflangja{19.2}{20.8}) {\iflangja{アンマーシャル}{unmarshal}};
\end{tikzpicture}
   (width ~117 mm, fits the 120 mm body)
\begin{tikzpicture}[x=1mm, y=1mm, font=\footnotesize\sffamily,
  >={Stealth[length=1.6mm]},
  cell/.style={draw, fill=white, minimum height=5mm, inner sep=0pt,
               font=\scriptsize\sffamily},
  var/.style={cell, minimum width=16mm},
  elem/.style={cell, minimum width=7mm},
  buf/.style={cell, minimum width=7mm, fill=ln-accent!15},
  vname/.style={font=\scriptsize\sffamily, anchor=east},
  hdr/.style={font=\scriptsize\sffamily\bfseries, text=ln-muted, anchor=base},
  lab/.style={font=\scriptsize\sffamily, text=ln-muted},
  addr/.style={font=\scriptsize\ttfamily, text=ln-muted, anchor=south west,
               inner sep=0.5pt},
  op/.style={font=\scriptsize\sffamily, anchor=south}]
  % ---------------- client address space
  \node[hdr] at (15,42) {\iflangja{クライアントのアドレス空間}{client address space}};
  \node[lab, anchor=west] at (0,36.5) {\iflangja{スタック}{stack}};
  \node[var] (cc) at (21,31) {3};
  \node[var] (cv) at (21,25) {\texttt{0x5a10}};
  \node[vname] at (12.5,31) {\texttt{count}};
  \node[vname] at (12.5,25) {\texttt{values}};
  \node[elem] (c1) at (11.5,9) {7};
  \node[elem] (c2) at (18.5,9) {$-2$};
  \node[elem] (c3) at (25.5,9) {300};
  \node[addr, anchor=north west] at (8,6.3) {0x5a10};
  \node[lab, anchor=east] at (7.6,9) {\iflangja{ヒープ}{heap}};
  \draw[->] ([xshift=-4mm]cv.south) -- (c1.north);
  % ---------------- message buffer
  \node[hdr] at (60,42) {\iflangja{メッセージバッファ}{message buffer}};
  \foreach \v [count=\i from 0] in {3,7,$-2$,300} {
    \node[buf] (b\i) at (49.5+7*\i,20) {\v};
    \node[font=\tiny\sffamily, text=ln-muted] at (46+7*\i,24.5) {\the\numexpr4*\i\relax};
  }
  \node[font=\tiny\sffamily, text=ln-muted] at (74,24.5) {16};
  \draw[ln-muted] (46.4,16.8) -- (46.4,15.8) -- (52.6,15.8) -- (52.6,16.8);
  \node[lab, anchor=north] at (49.5,15.3) {\texttt{count}};
  \draw[ln-muted] (53.4,16.8) -- (53.4,15.8) -- (73.6,15.8) -- (73.6,16.8);
  \node[lab, anchor=north] at (63.5,15.3)
    {\iflangja{要素}{elements}};
  % ---------------- server address space
  \node[hdr] at (103,42) {\iflangja{サーバのアドレス空間}{server address space}};
  \node[lab, anchor=west] at (88,36.5) {\iflangja{スタック}{stack}};
  \node[var] (sc) at (109,31) {3};
  \node[var] (sv) at (109,25) {\texttt{0x9c40}};
  \node[vname] at (100.5,31) {\texttt{count}};
  \node[vname] at (100.5,25) {\texttt{values}};
  \node[elem] (s1) at (99.5,9) {7};
  \node[elem] (s2) at (106.5,9) {$-2$};
  \node[elem] (s3) at (113.5,9) {300};
  \node[addr, anchor=north west] at (96,6.3) {0x9c40};
  \node[lab, anchor=east] at (95.6,9) {\iflangja{ヒープ}{heap}};
  \draw[->] ([xshift=-4mm]sv.south) -- (s1.north);
  % ---------------- operations
  % Japanese: below the arrows; above them the wider "アンマーシャル"
  % runs into the buffer offset "16".
  \draw[->, thick, ln-accent] (31,20) -- (44.5,20);
  \node[op, anchor=\iflangja{north}{south}] at (37.8,\iflangja{19.2}{20.8}) {\iflangja{マーシャル}{marshal}};
  \draw[->, thick, ln-accent] (75.5,20) -- (87,20);
  \node[op, anchor=\iflangja{north}{south}] at (\iflangja{82.3}{81.3},\iflangja{19.2}{20.8}) {\iflangja{アンマーシャル}{unmarshal}};
\end{tikzpicture}
   (width ~117 mm, fits the 120 mm body)
\begin{tikzpicture}[x=1mm, y=1mm, font=\footnotesize\sffamily,
  >={Stealth[length=1.6mm]},
  cell/.style={draw, fill=white, minimum height=5mm, inner sep=0pt,
               font=\scriptsize\sffamily},
  var/.style={cell, minimum width=16mm},
  elem/.style={cell, minimum width=7mm},
  buf/.style={cell, minimum width=7mm, fill=ln-accent!15},
  vname/.style={font=\scriptsize\sffamily, anchor=east},
  hdr/.style={font=\scriptsize\sffamily\bfseries, text=ln-muted, anchor=base},
  lab/.style={font=\scriptsize\sffamily, text=ln-muted},
  addr/.style={font=\scriptsize\ttfamily, text=ln-muted, anchor=south west,
               inner sep=0.5pt},
  op/.style={font=\scriptsize\sffamily, anchor=south}]
  % ---------------- client address space
  \node[hdr] at (15,42) {\iflangja{クライアントのアドレス空間}{client address space}};
  \node[lab, anchor=west] at (0,36.5) {\iflangja{スタック}{stack}};
  \node[var] (cc) at (21,31) {3};
  \node[var] (cv) at (21,25) {\texttt{0x5a10}};
  \node[vname] at (12.5,31) {\texttt{count}};
  \node[vname] at (12.5,25) {\texttt{values}};
  \node[elem] (c1) at (11.5,9) {7};
  \node[elem] (c2) at (18.5,9) {$-2$};
  \node[elem] (c3) at (25.5,9) {300};
  \node[addr, anchor=north west] at (8,6.3) {0x5a10};
  \node[lab, anchor=east] at (7.6,9) {\iflangja{ヒープ}{heap}};
  \draw[->] ([xshift=-4mm]cv.south) -- (c1.north);
  % ---------------- message buffer
  \node[hdr] at (60,42) {\iflangja{メッセージバッファ}{message buffer}};
  \foreach \v [count=\i from 0] in {3,7,$-2$,300} {
    \node[buf] (b\i) at (49.5+7*\i,20) {\v};
    \node[font=\tiny\sffamily, text=ln-muted] at (46+7*\i,24.5) {\the\numexpr4*\i\relax};
  }
  \node[font=\tiny\sffamily, text=ln-muted] at (74,24.5) {16};
  \draw[ln-muted] (46.4,16.8) -- (46.4,15.8) -- (52.6,15.8) -- (52.6,16.8);
  \node[lab, anchor=north] at (49.5,15.3) {\texttt{count}};
  \draw[ln-muted] (53.4,16.8) -- (53.4,15.8) -- (73.6,15.8) -- (73.6,16.8);
  \node[lab, anchor=north] at (63.5,15.3)
    {\iflangja{要素}{elements}};
  % ---------------- server address space
  \node[hdr] at (103,42) {\iflangja{サーバのアドレス空間}{server address space}};
  \node[lab, anchor=west] at (88,36.5) {\iflangja{スタック}{stack}};
  \node[var] (sc) at (109,31) {3};
  \node[var] (sv) at (109,25) {\texttt{0x9c40}};
  \node[vname] at (100.5,31) {\texttt{count}};
  \node[vname] at (100.5,25) {\texttt{values}};
  \node[elem] (s1) at (99.5,9) {7};
  \node[elem] (s2) at (106.5,9) {$-2$};
  \node[elem] (s3) at (113.5,9) {300};
  \node[addr, anchor=north west] at (96,6.3) {0x9c40};
  \node[lab, anchor=east] at (95.6,9) {\iflangja{ヒープ}{heap}};
  \draw[->] ([xshift=-4mm]sv.south) -- (s1.north);
  % ---------------- operations
  % Japanese: below the arrows; above them the wider "アンマーシャル"
  % runs into the buffer offset "16".
  \draw[->, thick, ln-accent] (31,20) -- (44.5,20);
  \node[op, anchor=\iflangja{north}{south}] at (37.8,\iflangja{19.2}{20.8}) {\iflangja{マーシャル}{marshal}};
  \draw[->, thick, ln-accent] (75.5,20) -- (87,20);
  \node[op, anchor=\iflangja{north}{south}] at (\iflangja{82.3}{81.3},\iflangja{19.2}{20.8}) {\iflangja{アンマーシャル}{unmarshal}};
\end{tikzpicture}
   (width ~117 mm, fits the 120 mm body)
\begin{tikzpicture}[x=1mm, y=1mm, font=\footnotesize\sffamily,
  >={Stealth[length=1.6mm]},
  cell/.style={draw, fill=white, minimum height=5mm, inner sep=0pt,
               font=\scriptsize\sffamily},
  var/.style={cell, minimum width=16mm},
  elem/.style={cell, minimum width=7mm},
  buf/.style={cell, minimum width=7mm, fill=ln-accent!15},
  vname/.style={font=\scriptsize\sffamily, anchor=east},
  hdr/.style={font=\scriptsize\sffamily\bfseries, text=ln-muted, anchor=base},
  lab/.style={font=\scriptsize\sffamily, text=ln-muted},
  addr/.style={font=\scriptsize\ttfamily, text=ln-muted, anchor=south west,
               inner sep=0.5pt},
  op/.style={font=\scriptsize\sffamily, anchor=south}]
  % ---------------- client address space
  \node[hdr] at (15,42) {\iflangja{クライアントのアドレス空間}{client address space}};
  \node[lab, anchor=west] at (0,36.5) {\iflangja{スタック}{stack}};
  \node[var] (cc) at (21,31) {3};
  \node[var] (cv) at (21,25) {\texttt{0x5a10}};
  \node[vname] at (12.5,31) {\texttt{count}};
  \node[vname] at (12.5,25) {\texttt{values}};
  \node[elem] (c1) at (11.5,9) {7};
  \node[elem] (c2) at (18.5,9) {$-2$};
  \node[elem] (c3) at (25.5,9) {300};
  \node[addr, anchor=north west] at (8,6.3) {0x5a10};
  \node[lab, anchor=east] at (7.6,9) {\iflangja{ヒープ}{heap}};
  \draw[->] ([xshift=-4mm]cv.south) -- (c1.north);
  % ---------------- message buffer
  \node[hdr] at (60,42) {\iflangja{メッセージバッファ}{message buffer}};
  \foreach \v [count=\i from 0] in {3,7,$-2$,300} {
    \node[buf] (b\i) at (49.5+7*\i,20) {\v};
    \node[font=\tiny\sffamily, text=ln-muted] at (46+7*\i,24.5) {\the\numexpr4*\i\relax};
  }
  \node[font=\tiny\sffamily, text=ln-muted] at (74,24.5) {16};
  \draw[ln-muted] (46.4,16.8) -- (46.4,15.8) -- (52.6,15.8) -- (52.6,16.8);
  \node[lab, anchor=north] at (49.5,15.3) {\texttt{count}};
  \draw[ln-muted] (53.4,16.8) -- (53.4,15.8) -- (73.6,15.8) -- (73.6,16.8);
  \node[lab, anchor=north] at (63.5,15.3)
    {\iflangja{要素}{elements}};
  % ---------------- server address space
  \node[hdr] at (103,42) {\iflangja{サーバのアドレス空間}{server address space}};
  \node[lab, anchor=west] at (88,36.5) {\iflangja{スタック}{stack}};
  \node[var] (sc) at (109,31) {3};
  \node[var] (sv) at (109,25) {\texttt{0x9c40}};
  \node[vname] at (100.5,31) {\texttt{count}};
  \node[vname] at (100.5,25) {\texttt{values}};
  \node[elem] (s1) at (99.5,9) {7};
  \node[elem] (s2) at (106.5,9) {$-2$};
  \node[elem] (s3) at (113.5,9) {300};
  \node[addr, anchor=north west] at (96,6.3) {0x9c40};
  \node[lab, anchor=east] at (95.6,9) {\iflangja{ヒープ}{heap}};
  \draw[->] ([xshift=-4mm]sv.south) -- (s1.north);
  % ---------------- operations
  % Japanese: below the arrows; above them the wider "アンマーシャル"
  % runs into the buffer offset "16".
  \draw[->, thick, ln-accent] (31,20) -- (44.5,20);
  \node[op, anchor=\iflangja{north}{south}] at (37.8,\iflangja{19.2}{20.8}) {\iflangja{マーシャル}{marshal}};
  \draw[->, thick, ln-accent] (75.5,20) -- (87,20);
  \node[op, anchor=\iflangja{north}{south}] at (\iflangja{82.3}{81.3},\iflangja{19.2}{20.8}) {\iflangja{アンマーシャル}{unmarshal}};
\end{tikzpicture}
